\begin{webpage}{calculus/limits/epsilon-delta-introduction}{Epsilon-Delta Definition: An Introduction}{lesson}
\chapter{Making Limits Precise}

\section{Why the Informal Definition Needs Numbers}

\lessonobjective{
Translate ``arbitrarily close'' into inequalities; understand the roles of \(\eps\) and \(\delta\); verify simple formal limit proofs.
}

The informal definition says that \(f(x)\) can be made as close as desired to \(L\) by taking \(x\) sufficiently close to \(a\). The phrase ``as close as desired'' must be made numerical.

\begin{bgconcept}
Suppose a customer says, ``Make the board close to 10 feet long.'' A builder needs a tolerance: within one inch? one millimeter? The symbol \(\eps\) names the allowed output error. The symbol \(\delta\) tells how close the input must be to guarantee that error.
\end{bgconcept}

The distance between \(x\) and \(a\) is
\[
  |x-a|.
\]
The distance between \(f(x)\) and \(L\) is
\[
  |f(x)-L|.
\]

\begin{bgdefinition}[title={Formal \(\eps\)-\(\delta\) Definition}]
We say
\[
  \lim_{x\to a}f(x)=L
\]
if
\[
  \boxed{\forall\eps>0,\ \exists\delta>0\text{ such that }
  0<|x-a|<\delta\implies|f(x)-L|<\eps.}
\]
\end{bgdefinition}

In words:

\begin{quote}
For every positive output tolerance \(\eps\), there is a positive input tolerance \(\delta\) such that every allowed input within \(\delta\) of \(a\), except \(a\) itself, produces an output within \(\eps\) of \(L\).
\end{quote}

\begin{center}
\begin{tabularx}{0.9\textwidth}{@{}p{0.24\textwidth}X@{}}
\toprule
\textbf{Expression} & \textbf{Meaning}\\
\midrule
\(|x-a|<\delta\) & The input lies in the horizontal interval \((a-\delta,a+\delta)\).\\
\(|f(x)-L|<\eps\) & The output lies in the vertical band \((L-\eps,L+\eps)\).\\
\(0<|x-a|\) & The target input itself is excluded.\\
\(\forall\eps>0\) & The argument must work for every requested accuracy.\\
\(\exists\delta>0\) & We may choose an input tolerance depending on \(\eps\).\\
\bottomrule
\end{tabularx}
\end{center}
\webnext{calculus/limits/epsilon-delta-linear}{Epsilon-Delta Proofs for Linear Functions}
\end{webpage}

\begin{webpage}{calculus/limits/epsilon-delta-linear}{Epsilon-Delta Proofs for Linear Functions}{lesson}
\section{Linear Functions: The Cleanest Proofs}

\lessonobjective{
Construct a \(\delta\) directly from a requested \(\eps\) for a linear function.
}

\begin{bgguided}[title={\(f(x)=2x\)}]
Prove
\[
  \lim_{x\to3}2x=6.
\]

\begin{bgsolution}
We want the output error to be less than \(\eps\):
\[
  |2x-6|<\eps.
\]
Factor:
\[
  |2x-6|=2|x-3|.
\]
So it is enough to require
\[
  2|x-3|<\eps.
\]
Divide by \(2\):
\[
  |x-3|<\frac\eps2.
\]
Choose
\[
  \boxed{\delta=\frac\eps2}.
\]
Then if \(0<|x-3|<\delta\),
\[
  |2x-6|=2|x-3|<2\delta=2\left(\frac\eps2\right)=\eps.
\]
Therefore,
\[
  \lim_{x\to3}2x=6.
\]
\end{bgsolution}
\end{bgguided}

\begin{bgcheck}{epsilon-flow-01}{expression}{epsilon/2}[Check your understanding]
For \(f(x)=2x\) at \(x=3\), choose \(\delta\) in terms of \(\varepsilon\).
\begin{bghint}
\(|2x-6|=2|x-3|\).
\end{bghint}
\begin{bgreveal}
Choose \(\delta=\varepsilon/2\).
\end{bgreveal}
\end{bgcheck}


This proof has two phases.

\begin{description}
  \item[Discovery phase.] Start from the desired output inequality and work backward to guess a useful \(\delta\).
  \item[Proof phase.] State the chosen \(\delta\), assume the input inequality, and work forward to the output inequality.
\end{description}

\begin{bgexample}[title={A general linear function}]
Prove
\[
  \lim_{x\to2}(3x+1)=7.
\]

\begin{bgsolution}
\textbf{Discovery.} We need
\[
  |(3x+1)-7|<\eps.
\]
Simplify:
\[
  |3x-6|=3|x-2|.
\]
Thus, requiring
\[
  |x-2|<\frac\eps3
\]
is enough.

\textbf{Proof.} Let \(\eps>0\). Choose
\[
  \delta=\frac\eps3.
\]
If
\[
  0<|x-2|<\delta,
\]
then
\[
  |(3x+1)-7|=3|x-2|<3\delta=\eps.
\]
Therefore,
\[
  \boxed{\lim_{x\to2}(3x+1)=7}.
\]
\end{bgsolution}
\end{bgexample}

\begin{bgsummary}
For \(f(x)=mx+b\), a standard choice is
\[
  \delta=\frac\eps{|m|}
\]
when \(m\ne0\). The output changes \(|m|\) times as much as the input, so the input tolerance must be \(|m|\) times smaller.
\end{bgsummary}
\webnext{calculus/limits/epsilon-delta-quadratic}{Epsilon-Delta Proofs for Quadratic Functions}
\end{webpage}

\begin{webpage}{calculus/limits/epsilon-delta-quadratic}{Epsilon-Delta Proofs for Quadratic Functions}{lesson}
\section{Nonlinear Functions: Controlling an Extra Factor}

\lessonobjective{
Use a preliminary restriction such as \(|x-a|<1\) to bound a factor that depends on \(x\); choose \(\delta\) as the minimum of two requirements.
}

For a quadratic, factoring the output error creates an extra factor involving \(x\).

\begin{bgexample}[title={Prove \(\lim_{x\to2}x^2=4\)}]
\begin{bgsolution}
We need to make
\[
  |x^2-4|
\]
small. Factor:
\[
  |x^2-4|=|x-2||x+2|.
\]
The first factor is controlled by \(\delta\), but \(|x+2|\) still depends on \(x\). First impose the simple restriction
\[
  |x-2|<1.
\]
Then
\[
  1<x<3,
\]
so
\[
  3<x+2<5,
\]
and therefore
\[
  |x+2|<5.
\]
Now
\[
  |x^2-4|=|x-2||x+2|<5|x-2|.
\]
To make this less than \(\eps\), require
\[
  |x-2|<\frac\eps5.
\]
We need both restrictions, so choose
\[
  \boxed{\delta=\min\left\{1,\frac\eps5\right\}}.
\]

Now write the forward proof. Let \(\eps>0\), and choose the \(\delta\) above. If
\[
  0<|x-2|<\delta,
\]
then \(|x-2|<1\), so \(|x+2|<5\). Also, \(|x-2|<\eps/5\). Hence
\[
  |x^2-4|
  =|x-2||x+2|
  <5\left(\frac\eps5\right)
  =\eps.
\]
Therefore,
\[
  \boxed{\lim_{x\to2}x^2=4}.
\]
\end{bgsolution}
\end{bgexample}

\begin{bgconcept}
The minimum notation means, ``Use whichever restriction is smaller.'' If \(\eps/5\) is tiny, use it. If \(\eps/5\) is larger than \(1\), keep the local bound \(1\). Both jobs must be done at once.
\end{bgconcept}

\begin{bgexample}[title={A shifted square}]
Prove
\[
  \lim_{x\to1}x^2=1.
\]

\begin{bgsolution}
Factor the output error:
\[
  |x^2-1|=|x-1||x+1|.
\]
If \(|x-1|<1\), then \(0<x<2\), so \(|x+1|<3\). Therefore,
\[
  |x^2-1|<3|x-1|.
\]
Choose
\[
  \boxed{\delta=\min\left\{1,\frac\eps3\right\}}.
\]
Then
\[
  |x^2-1|<3\delta\le\eps.
\]
\end{bgsolution}
\end{bgexample}
\webnext{calculus/limits/epsilon-delta-graph}{Reading Epsilon and Delta From a Graph}
\end{webpage}

\begin{webpage}{calculus/limits/epsilon-delta-graph}{Reading Epsilon and Delta From a Graph}{lesson}
\section{Reading the Definition From a Graph}

\begin{figure}[htbp]
\centering
\begin{tikzpicture}
\begin{axis}[
  width=0.84\textwidth,height=0.50\textwidth,
  xmin=-0.5,xmax=4.5,ymin=-0.5,ymax=8.5,
  axis lines=middle,xlabel={$x$},ylabel={$f(x)$},
  grid=major,major grid style={draw=BGLine}
]
\addplot[BGBlue,very thick,domain=-0.3:4.2,samples=140]{x^2/2+1};
\addplot[draw=none,name path=upper,domain=-0.5:4.5]{5.5};
\addplot[draw=none,name path=lower,domain=-0.5:4.5]{3.5};
\addplot[BGPaleGold] fill between[of=upper and lower];
\draw[BGGold,thick,dashed] (axis cs:-0.5,5.5)--(axis cs:4.5,5.5);
\draw[BGGold,thick,dashed] (axis cs:-0.5,3.5)--(axis cs:4.5,3.5);
\draw[BGGreen,thick,dashed] (axis cs:1.75,-0.5)--(axis cs:1.75,8.5);
\draw[BGGreen,thick,dashed] (axis cs:2.25,-0.5)--(axis cs:2.25,8.5);
\node[anchor=west,font=\small,BGGold] at (axis cs:3.65,5.5) {$L+\eps$};
\node[anchor=west,font=\small,BGGold] at (axis cs:3.65,3.5) {$L-\eps$};
\node[anchor=north,font=\small,BGGreen] at (axis cs:1.75,-0.2) {$a-\delta$};
\node[anchor=north,font=\small,BGGreen] at (axis cs:2.25,-0.2) {$a+\delta$};
\end{axis}
\end{tikzpicture}
\caption{An \(\eps\)-band around \(L\) and a corresponding \(\delta\)-window around \(a\). A valid \(\delta\) keeps the nearby graph inside the band.}
\label{fig:epsilon-delta}
\end{figure}

\begin{bggraphspec}
Use a continuous nonlinear function such as \(f(x)=x^2/2+1\) near \(a=2\), where \(L=3\). Provide an adjustable \(\eps\)-band and show the largest symmetric \(\delta\)-window whose graph segment remains inside the band. Display \(|x-a|<\delta\) and \(|f(x)-L|<\eps\) simultaneously.
\end{bggraphspec}
\webnext{calculus/limits/disprove-limit-epsilon-delta}{How to Disprove a Claimed Limit}
\end{webpage}

\begin{webpage}{calculus/limits/disprove-limit-epsilon-delta}{How to Disprove a Claimed Limit}{lesson}
\section{Proving a Limit Is Not a Claimed Number}

To show
\[
  \lim_{x\to a}f(x)\ne L,
\]
it is enough to find one positive tolerance \(\eps_0\) for which no possible \(\delta\) works.

\begin{bgexample}[title={A jump fails the formal definition}]
Let
\[
  f(x)=
  \begin{cases}
    0,&x<0,\\
    1,&x\ge0.
  \end{cases}
\]
Show that \(\lim_{x\to0}f(x)\) does not equal \(1/2\).

\begin{bgsolution}
Choose
\[
  \eps_0=\frac14.
\]
The band around \(1/2\) is
\[
  \left(\frac14,\frac34\right).
\]
But every left-side input near zero has output \(0\), which lies outside this band, and every right-side input has output \(1\), also outside the band.

No matter how small \(\delta>0\) is chosen, the point \(x=-\delta/2\) satisfies
\[
  0<|x|<\delta,
\]
but
\[
  \left|f(x)-\frac12\right|
  =\left|0-\frac12\right|
  =\frac12
  >\frac14.
\]
Therefore the claimed limit cannot be \(1/2\). In fact, the left and right limits disagree, so no two-sided limit exists.
\end{bgsolution}
\end{bgexample}
\webnext{calculus/limits/unit-review}{Limits and Continuity Unit Review}
\end{webpage}

\begin{webpage}{calculus/limits/formal-infinite-limits}{Formal Definition of an Infinite Limit}{extension}
\section{Formal Infinite Limits: Optional Extension}

The statement
\[
  \lim_{x\to a}f(x)=+\infty
\]
can be formalized as follows:

\begin{bgdefinition}[title={Formal Positive Infinite Limit}]
For every \(M>0\), there exists \(\delta>0\) such that
\[
  0<|x-a|<\delta\implies f(x)>M.
\]
\end{bgdefinition}

Instead of asking for an output band around a finite \(L\), the challenger names an arbitrarily high threshold \(M\). The proof must force the function above it.

\begin{bgguided}[title={\(1/x^2\to+\infty\)}]
Show formally that
\[
  \lim_{x\to0}\frac1{x^2}=+\infty.
\]

\begin{bgsolution}
Let \(M>0\). We want
\[
  \frac1{x^2}>M.
\]
This is equivalent to
\[
  x^2<\frac1M,
\]
which is guaranteed by
\[
  |x|<\frac1{\sqrt M}.
\]
Choose
\[
  \delta=\frac1{\sqrt M}.
\]
Then \(0<|x|<\delta\) implies
\[
  x^2<\delta^2=\frac1M,
\]
so
\[
  \frac1{x^2}>M.
\]
\end{bgsolution}
\end{bgguided}
\webnext{calculus/limits/unit-review}{Limits and Continuity Unit Review}
\end{webpage}

\begin{webpage}{calculus/limits/epsilon-delta-review}{Epsilon-Delta Review}{review}
\section*{Chapter 6 Summary}
\addcontentsline{toc}{section}{Chapter 6 Summary}

\begin{bgsummary}
\begin{itemize}
  \item \(\eps\) is an allowed output error; \(\delta\) is an input tolerance that guarantees it.
  \item Formal proofs work backward to discover \(\delta\) and forward to verify it.
  \item Linear proofs usually use \(\delta=\eps/|m|\).
  \item Nonlinear proofs often need a preliminary local bound and a minimum of two restrictions.
  \item To disprove a claimed limit, one bad \(\eps_0\) is enough.
\end{itemize}
\end{bgsummary}
\webnext{calculus/limits/unit-review}{Limits and Continuity Unit Review}
\end{webpage}

\begin{webpage}{calculus/limits/epsilon-delta-concept-quiz}{Epsilon-Delta Concept Quiz}{quiz}
\section*{Chapter 6 Concept Quiz}
\addcontentsline{toc}{section}{Chapter 6 Concept Quiz}

Use this as a short retrieval check after finishing the chapter. On the website, each item should accept an answer before revealing the hint and worked explanation.

\begin{bgcheck}{epsilon-linear-q1}{expression}{epsilon/3}[Concept check]
For \(f(x)=3x+1\) at \(x=2\), give a working choice of \(\delta\) in terms of \(\varepsilon\).
\begin{bghint}
\(|f(x)-7|=3|x-2|\).
\end{bghint}
\begin{bgreveal}
Choose \(\delta=\varepsilon/3\).
\end{bgreveal}
\end{bgcheck}

\begin{bgcheck}{epsilon-linear-q2}{expression}{epsilon/5}[Concept check]
For \(f(x)=5x\) at \(x=1\), give a working \(\delta\).
\begin{bghint}
\(|5x-5|=5|x-1|\).
\end{bghint}
\begin{bgreveal}
Choose \(\delta=\varepsilon/5\).
\end{bgreveal}
\end{bgcheck}

\begin{bgcheck}{epsilon-quadratic-q1}{expression}{min(1,epsilon/5)}[Concept check]
For the standard proof that \(x^2\to4\) as \(x\to2\), enter the usual safe choice of \(\delta\).
\begin{bghint}
First require \(|x-2|<1\), which makes \(|x+2|<5\).
\end{bghint}
\begin{bgreveal}
Choose \(\delta=\min\{1,\varepsilon/5\}\).
\end{bgreveal}
\end{bgcheck}

\begin{bgcheck}{formal-reading-q1}{choice}{output}[Concept check]
In \(|f(x)-L|<\varepsilon\), does \(\varepsilon\) control input distance or output distance?
\begin{bghint}
Look at which quantities are being compared.
\end{bghint}
\begin{bgreveal}
It controls output distance.
\end{bgreveal}
\end{bgcheck}

\begin{bgsummary}[title={Quiz use on the website}]
This page should report completion by concept, not merely a total score. A wrong answer should link back to the exact lesson that teaches the missing method. The same questions may appear in randomized numerical variants, but the canonical explanation should remain stable.
\end{bgsummary}
\webnext{calculus/limits/unit-review}{Limits and Continuity Unit Review}
\end{webpage}

\begin{webpage}{calculus/limits/epsilon-delta-practice}{Epsilon-Delta Practice Problems}{practice}
\section*{Chapter 6 Exercises}
\addcontentsline{toc}{section}{Chapter 6 Exercises}

\begin{bgexercises}[label=\textbf{6.\arabic*}]
  \item In your own words, explain the roles of \(\eps\) and \(\delta\).
  \item Rewrite \(|x-3|<0.2\) as an ordinary interval.
  \item Rewrite \(|f(x)-5|<0.1\) as a vertical output interval.
  \item For \(f(x)=4x\), find a \(\delta\) in terms of \(\eps\) that proves \(\lim_{x\to2}f(x)=8\).
  \item Prove \(\lim_{x\to1}(5x-2)=3\).
  \item Prove \(\lim_{x\to-2}(3x+4)=-2\).
  \item Prove \(\lim_{x\to4}(x/2)=2\).
  \item Prove \(\lim_{x\to1}x^2=1\) using \(\delta=\min\{1,\eps/3\}\).
  \item Prove \(\lim_{x\to3}x^2=9\) by first bounding \(|x+3|\).
  \item Find a suitable \(\delta\) to prove \(\lim_{x\to2}(x^2+x)=6\).
  \item Explain why a proof may choose a smaller \(\delta\) than necessary.
  \item Explain why \(\delta\) may depend on \(\eps\), but not on the particular \(x\) chosen later.
  \item Use \(\eps_0=1/4\) to show the step function in the chapter does not approach \(1/2\) at zero.
  \item Give a formal \(M\)-\(\delta\) proof that \(1/x^2\to+\infty\) as \(x\to0\).
  \item Give a formal \(N\)-style proof that \(1/x\to0\) as \(x\to\infty\).
  \item Identify the first incorrect step in the claim: ``Choose \(\delta=\eps\). Then \(|x^2-4|<\eps\) whenever \(|x-2|<\delta\).''
  \item Why is \(0<|x-a|\) included in the formal finite-limit definition?
  \item Explain how the graphical \(\eps\)-band and \(\delta\)-window represent the quantified definition.
\end{bgexercises}

Answers begin in \cref{app:chapter-answers}.
\webnext{calculus/limits/unit-review}{Limits and Continuity Unit Review}
\end{webpage}
